\section{Channels and closure under selected laws}\label{sec:channels}

The selected length is determined by the original point in the suspension.
It can therefore change the joint law of other coordinates, even when
their unconditioned law is a product.  We organize those coordinates in
three groups.  Independence of the original point from any two groups
will control the effect of selecting a length on all three groups.
This channel construction and its closure argument adapt
\cite[Section~5]{FourfoldCompanion}; the proof below permits every fixed
finite set of slots $I=\{1,\ldots,n\}$.

\begin{definition}[Channel]\label{def:channel}
A \emph{channel} is a flow $Y$ over the suspension $C$, with factor map
$\pi_Y:Y\to C$, together with a finite list of probability-preserving
maps $D_\gamma:Y\to C$, called \emph{outputs}, and positive rational
numbers $q_\gamma$, such that
\[
 D_\gamma(S_Y^t y)=S_C^{q_\gamma t}D_\gamma(y)
\]
for each fixed $t$, modulo null sets.  Partition the output indices into
three groups, numbered $1,2,3$; empty groups are allowed.  Set
\[
 x_0=\xi\circ\pi_Y,\qquad x_\gamma=\xi\circ D_\gamma,
 \qquad d_Y=(x_\gamma)_\gamma,
\]
and let $d_{Y,H}$ be the sublist from groups in
$H\subseteq\{1,2,3\}$.  Let $V$ collect the phase coordinates of
$\pi_Y$ and all the $D_\gamma$.  Under the measure $m_Y$ require
\begin{equation}\label{eq:channel-laws}
 \begin{split}
  \law(d_Y\mid V)&=\bigotimes_\gamma\mu,\\
  \law(x_0,d_{Y,H}\mid V)&=
       \mu\otimes\bigotimes_{\gamma\text{ in groups }H}\mu,
       \qquad H\subsetneq\{1,2,3\}.
 \end{split}
\end{equation}
An empty spatial list has the one-point law.
\end{definition}

Thus the spatial outputs have a product law independent of the phases.
The original spatial coordinate has a product law with the outputs in
any proper collection of groups, but no independence from the complete
output list is required.  The suspension $C$, with no outputs, is a
channel.

\begin{lemma}[Conditioning and extensions]\label{lem:channel-conditioning}
For a channel $Y$, both identities in \eqref{eq:channel-laws} hold with
conditioning on its maximal measure-distal factor $W_Y$ in place of
$V$.  Every flow extension of $Y$ over $C$ is a channel with the pulled
back output data, and the same strengthening holds there.
\end{lemma}

\begin{proof}
Choose a positive integer $h$ such that $hq_\gamma$ is an integer for
each output.  The phase vector $V$ is a factor fixed by time $h$, and
its time-one action is distal, so $V$ belongs to $W_Y$.  For any spatial
tuple appearing in \eqref{eq:channel-laws}, its joint factor with $V$
has product law over $V$.  At time $h$ the spatial coordinates evolve
by a product of positive integer powers of $T$.  This product is mixing:
the assertion follows first for product tests from mixing of each
power, and then for all bounded tests by approximation.  Consequently
this factor is relatively weakly mixing over $V$ for time $h$.

On the other hand, $W_Y$ is distal for time $h$ and relatively distal
over $V$.  Relative disjointness, in the form of
\cite[Proposition~2.2]{FourfoldCompanion}, makes the spatial factor and
$W_Y$ relatively independent over $V$.  Its conditional spatial law
given $W_Y$ is therefore the product law in
\eqref{eq:channel-laws}.  Pullback to an extension preserves the joint
laws of the original data, proving the channel identities there; the
same argument gives their stronger conditioning.
\end{proof}

We now prove the closure property that will permit repeated extensions.
Recall from Lemma~\ref{lem:normalized-joining} that
$\mathsf M_b(Y)$ has state space $Y^I$, measure $\lambda_Y$, and flow
$(y_i)_i\mapsto(S_Y^{it/b}y_i)_i$.

\begin{proposition}[Channel closure]\label{prop:channel-closure}
Let $Y$ be a channel and $b\in I$.  On $\mathsf M_b(Y)$ take
$\pi_Y(y_b)$ as original factor and take every $D_\gamma(y_i)$,
$i\in I$, as an output, in the old group of $\gamma$.
These data make $\mathsf M_b(Y)$ a channel.
\end{proposition}

\begin{proof}
Each new output is probability-preserving and has positive rational
speed $iq_\gamma/b$.  It remains to prove the two conditional laws.
Test them with products of bounded spatial functions and products of
bounded phase functions.  Such tests determine the laws.  Expanding
each spatial function into its mean and centered part, and scaling,
reduces the proof to real spatial functions that are either $1$ or
centered and bounded in absolute value by $1$.  We may also bound every
phase function by $1$.  Whenever at least one spatial factor is
centered, the tested integral must be zero.  A test using the new
original coordinate uses outputs from only a proper set of groups.

\paragraph{Time intervals and offsets.}
Use the selected-law formula \eqref{eq:selected-law} and condition on
the old phase vector $V$ under $m_Y$.  Fix an integer $h>0$ clearing
all output-speed denominators and write
$t=hl+s$, where $l\ge0$ is an integer and $0\le s<h$.
For fixed $V,s$, all phase weights in the test are independent of $l$.
The spatial function for output $\gamma$ at slot $i$ is evaluated at
\begin{equation}\label{eq:channel-times}
 T^{e_\gamma i l+a_{\gamma,i}}x_\gamma,
 \qquad e_\gamma=hq_\gamma\in\N.
\end{equation}
The offsets $a_{\gamma,i}$ depend on $V,s$ and range over a fixed
finite set, by the suspension formula.  If the original-factor test
is present, its argument is $T^{hbl+a_0}x_0$, with the same finite-offset
property.  Equivariance and Fubini justify these formulas in the time
integrals.

On an event where the selected length is $N=c2^k$, $c\in F$ and
$k\ge a$, put $m=\lfloor c2^k/h\rfloor$.  Replacing the normalized
time integral by the average over the $m$ complete intervals, followed
by integration against $ds/h$, has error at most $2h2^{-a}$.
For all sufficiently large $a$ we have $m\ge1$.

\paragraph{Tests using a proper set of groups.}
Given $V$, the relevant old output coordinates have product law and
are independent of $x_0$, by \eqref{eq:channel-laws}.  The selector is
determined by $V,x_0$.  Suppose first that some output occurrence is
centered.  For fixed $V,s,x_0$, integrate the spatial output product
at index $l$.  The expectation factors over the old output indices
$\gamma$.  For an index having a centered occurrence, the times
\eqref{eq:channel-times} in distinct slots separate as $l\to\infty$.
Mixing of all orders makes its integral tend to the product of the
means, hence to zero.  A single centered occurrence already has zero
integral.  The convergence is uniform over the finitely many offset
patterns.  All other factors have absolute value at most $1$.
The Ces\`aro averages of these absolute expectations consequently tend
to zero uniformly over all sufficiently large lengths.  This applies
to the selected lengths as well.

If every output test is $1$, the centered factor must be the original
one.  Birkhoff's theorem for the ergodic power $T^{hb}$ gives convergence
to zero for $\mu$-almost every $x_0$, simultaneously for all the
finitely many offsets $a_0$.  Since the selected lengths tend to
infinity pointwise, they have the same limit.  The conditional law of
$x_0$ given $V$ is $\mu$, so bounded convergence completes this case.

\paragraph{Three groups and almost orthogonal vectors.}
It remains to treat the full output list, with no original-factor
spatial test.  If one group has no centered occurrence, every spatial
test in that group is $1$, and the preceding argument applies; the
phase tests of that group were allowed throughout.  Thus suppose each
group contains a centered occurrence.  For fixed $V,s$, write its
spatial product at index $l$ as
\[
 u_l\in L^2(\Omega_1),\qquad
 v_l\in L^2(\Omega_2),\qquad
 w_l\in L^2(\Omega_3),
\]
where $\Omega_j$ carries the product measure with one copy of $\mu$
for every old output in group $j$.  These are real unit-ball vectors,
indeed functions bounded by $1$.  Given $V$, the old spatial output
law is the product of these three spaces.  The vectors depend on
$V,s$ only through the finite offset patterns.

For every $\eta>0$, after omitting finitely many initial indices, the
families $(u_l)$ and $(v_l)$ each have absolute pair inner products at
most $\eta$ outside a graph of bounded maximum degree $L_0$.
All these bounds are uniform in the offset patterns.  To verify this,
in each of the two groups choose a centered occurrence at some output
index $\gamma$.  Its coordinate integral in a pair inner product is
small by mixing of all orders whenever the times in the two copies of
\eqref{eq:channel-times} are sufficiently separated.  The times within
each copy separate once $l,l'$ exceed a fixed threshold.  Failure of
separation between copies then requires one of finitely many inequalities
\[
 \abs{e_\gamma i l-e_\gamma i'l'+a'}\le R.
\]
For fixed $l$, each inequality permits only a bounded number of
integers $l'$, because $e_\gamma,i'>0$; the same holds with $l,l'$
interchanged.  Their union is the required bounded-degree graph.
The coordinate integrals for other output indices are bounded by $1$.
The separation threshold and $L_0$ may depend on $\eta$, but no rate
of mixing is needed.

\paragraph{Conditioning on a selected length.}
For fixed selector index $a$, partition into the events
$N_a=c2^k$, $c\in F$, $k\ge a$, choosing one representation when a
length has more than one.  Given $V$, write their probabilities as
$p_{k,c}(V)$.  On one such event of probability $p>0$, the conditional
spatial law of any two groups, and of each single group, is unchanged.
Indeed the event is determined by $V,x_0$, which is independent of any
proper set of groups by \eqref{eq:channel-laws}.  The full three-group
law may change.  Its density $\rho$ relative to the product law satisfies
\begin{equation}\label{eq:channel-density}
 0\le p\rho\le1,\qquad
 \int p\rho=p,\qquad \norm{p\rho}_2\le\sqrt p.
\end{equation}
The reason is that the joint measure of the outputs and the event is
dominated by the unconditioned product measure.  These statements hold
for almost every $V$; the length events are countable.

We choose the rank bound to meet two requirements: the packing
bound must be $o(m)$, and the reciprocal bounds must be summable over
dyadic length scales. For $m=\lfloor c2^k/h\rfloor$ set
\begin{equation}\label{eq:channel-dimension}
 d_m=\left\lceil(1+\log m)^{3/2}\right\rceil,
 \qquad z_l=u_l\otimes v_l\quad(0\le l<m).
\end{equation}
Let $\Pi_m$ be the orthogonal projection onto up to $d_m$ leading
eigenvectors of the positive covariance operator
$\sum_{l<m}z_lz_l^*$, taking its full range if its rank is smaller.
Fix $0<\varepsilon\le1$ and then fix
$0<\eta\le\varepsilon^2/10^4$ in the preceding graph bound.
Theorem~\ref{thm:packing}, proved in Appendix~\ref{sec:packing} from
\cite[Theorem~4.1]{FourfoldCompanion}, gives
\begin{equation}\label{eq:channel-projected-average}
 \frac1m\sum_{l<m}\norm{\Pi_m z_l}
       \le\varepsilon+o(1).
\end{equation}
Here the error tends to zero uniformly over the offset patterns.
Indeed, after the fixed initial indices are removed, the number of
projection norms at least $\varepsilon$ is at most
\[
 C(1+L_0)\varepsilon^{-2}
 \exp\!\bigl(C\varepsilon^{-2}\sqrt{d_m}\log(2+d_m)\bigr)=o(m).
\]
The removed indices contribute $o(1)$ to the average.  Importantly,
$\varepsilon,\eta$, and the resulting exceptional degree are fixed
before $m$ tends to infinity.

For the projected terms, Cauchy--Schwarz under the conditional event
law uses its unchanged two-group and one-group marginals.  Term by
term, including the event probability, it gives
\[
 \left|p\int (\Pi_mz_l)(\omega_1,\omega_2)
                   w_l(\omega_3)\rho(\omega)\,d\omega\right|
 \le p\norm{\Pi_mz_l}_2\norm{w_l}_2
 \le p\norm{\Pi_mz_l}_2.
\]
Thus their average contributes at most $p(\varepsilon+o(1))$.

For the residual vectors $z_l'=(1-\Pi_m)z_l$, let
$A=(\langle z_l',z_{l'}'\rangle)_{l,l'<m}$ and
$B=(\langle w_l,w_{l'}\rangle)_{l,l'<m}$.  The covariance trace is
at most $m$, so its remaining eigenvalues, and hence the nonzero
eigenvalues of $A$, are at most $m/d_m$.  If its rank was at most
$d_m$, the residual is zero.  Also $B$ is positive semidefinite with
trace at most $m$.  Consequently
\[
 \left\|\frac1m\sum_{l<m}z_l'\otimes w_l\right\|^2
   =\frac{\operatorname{tr}(AB)}{m^2}\le\frac1{d_m}.
\]
By \eqref{eq:channel-density}, the residual contribution, including
the event probability, is at most $\sqrt{p/d_m}$.
Multiplication by the fixed phase weight does not increase either bound.

\paragraph{Summing over lengths.}
Write $m(k,c)=\lfloor c2^k/h\rfloor$.  For fixed $h$ this is comparable
to $2^k$ for all sufficiently large $k$, uniformly in the finite set
$F$.  Since $\sum_{k,c}p_{k,c}(V)=1$, Cauchy--Schwarz yields
\begin{equation}\label{eq:channel-residual-sum}
 \sum_{k\ge a,\ c\in F}\sqrt{\frac{p_{k,c}(V)}{d_{m(k,c)}}}
 \le
 \left(\sum_{k\ge a,\ c\in F}\frac1{d_{m(k,c)}}\right)^{1/2}
 \longrightarrow0.
\end{equation}
The series converges because its summands are $O(k^{-3/2})$.
The projected contributions have total limit superior at most
$\varepsilon$, by \eqref{eq:channel-projected-average} and the same
probability identity.  These estimates are uniform in $V,s$ and may
be integrated.  There is no measurable-choice issue for the
projections: choose one for each of the countably many sizes and
finitely many offset patterns.  Letting $a$ tend to infinity in the
selected-law formula and then letting $\varepsilon$ tend to zero
proves the required vanishing.  Both laws in
\eqref{eq:channel-laws} follow.
\end{proof}

\begin{corollary}\label{cor:channel-copied-law}
For every channel $Y$, the copied spatial list
$(d_Y(y_i))_{i\in I}$ under $\lambda_Y$ has its full product law
conditional on $P=(\omega_Y(y_i))_{i\in I}$.
\end{corollary}

\begin{proof}
Apply Proposition~\ref{prop:channel-closure} and
Lemma~\ref{lem:channel-conditioning} to $\mathsf M_b(Y)$.
By Lemma~\ref{lem:normalized-joining}, the old parameter array $P$
belongs to its maximal distal factor.  Conditioning the resulting
constant product law down to $P$ proves the assertion.
\end{proof}
